\scanpage{152}\mspage{39}
regular map of $\widetilde{V}$ onto $V$ which is bi-regular between $\widetilde{V}-S$ and $V-\mathfrak{p}$. It is easy to prove the following propositions:\footnotemark[4]

i) $Q_{\mathfrak{p}}^{-1}$ induces an isomorphism: $\unsure{\mathfrak{M}|\mathscr{M}}(Q_{\mathfrak{p}}(V))\cong\mathfrak{M}(V)$.

ii) If $V$ is a non-singular algebraic surface in a projective space, then $Q_{\mathfrak{p}}(V)$ is also a non-singular algebraic surface in a projective space.

Moreover we have

\begin{quote}
CASTELNUOVO-ENRIQUES' THEOREM:\footnotemark[5] Let $\widetilde{V}$ be a non-singular algebraic surface in a projective space and let $S$ be a non-singular rational curve on $\widetilde{V}$ with $(S^{2})=-1$. Then there exists a non-singular algebraic surface $V$ in a projective space and a point $\mathfrak{p}$ on $V$ such that $Q_{\mathfrak{p}}(V)=\widetilde{V}$ and $Q_{\mathfrak{p}}(\mathfrak{p})=S$.
\end{quote}

The following proposition\footnotemark[6] is a corollary to this theorem:

iii) If $Q_{\mathfrak{p}}(V)$ is a non-singular algebraic surface in a projective space, then $V$ is also a non-singular algebraic surface in a projective space.
\footnotetext[4]{See Chow and Kodaira, loc. cit.}
\footnotetext[5]{G. Castelnuovo and F. Enriques, Sopra alcune questioni fondamentali nella teoria delle superficie algebriche, Annali di Math., Ser. III, Vol. 6 (1901), pp. 165-225, Chap. III. See also K. Kodaira, On Kähler varieties of restricted type, Annals of Math., Vol. 60 (1954), pp. 44-48.}
\footnotetext[6]{Chow and Kodaira, loc. cit., Lemma 3.}
% archive sheet number: 151
